% !Mode:: "TeX:UTF-8"
\chapter{基于威胁势场与多头注意力机制的多无人艇协同围捕博弈策略}
\label{cha:hunting}

\section{引言}
\label{sec:hunt_intro}
相较于单艇追逃博弈, 多无人艇协同围捕博弈不仅要求追击者逼近逃逸者, 还需要多艇之间
在空间上形成包围态势, 共同限制逃逸者的机动空间直至完成捕获. 在含障碍物的海洋环
境中, 协同围捕面临以下三方面挑战: 其一, 现有研究大多未给出严格的围捕成功数学判
据, 导致算法评价依据模糊; 其二, 围捕过程呈现明显的阶段性特征 (接近、成型、保持),
单一的奖励函数难以同时兼顾各阶段的学习需求, 常导致稀疏奖励与训练收敛缓慢; 其三,
现有工作多止步于数值仿真, 缺乏真实水面环境下的物理平台验证, 算法的工程可实现性
有待考察.

针对上述问题, 本章围绕多无人艇协同围捕博弈策略的设计与验证, 开展以下四方面工作:
(i) 在六自由度欠驱动运动学模型的基础上, 给出围捕成功的显式数学判据, 解决现有研
究中评价标准模糊的问题; (ii) 提出能量波峰 (Energy Crest, EC) 与能量波谷 (Energy
Trough, ET) 概念, 将威胁势场转化为反映接近-成型-保持三阶段的密集奖励机制; (iii)
在 MAPPO 算法的 Critic 网络中引入多头注意力机制, 形成 MHD-MAPPO 算法; (iv)
\textbf{搭建多无人艇物理实验平台, 在真实水面环境中完成 2 对 1 和 3 对 1 协同围捕
实物实验}, 从仿真到实物全面验证所提算法. 本章内容在作者已发表
的文献~\cite{wang2025cooperativehunting} 基础上整理扩展而成.

\section{复杂环境下多无人艇协同围捕问题}
\label{sec:hunt_problem}

\subsection{问题描述}
\label{subsec:hunt_scene}
考虑含 $N$ 艘追击者 $P_n\ (n=1,2,\ldots,N)$、$M$ 艘逃逸者 $E_m\ (m=1,2,\ldots,M)$
与 $U$ 个静态障碍物 $\mathrm{obs}_u\ (u=1,2,\ldots,U)$ 的复杂海洋博弈场景, 如
图~\ref{fig:hunt_scene} 所示. 追击者协同完成对逃逸者的合围任务, 同时避开障碍物
及彼此; 逃逸者则根据周围态势采取合适的逃逸策略. 相关变量定义如下:
$\mathbf{v}_i$ 为第 $i$ 艘无人艇的速度, $\varphi_i\in[0, 2\pi)$ 为艏向角,
$d_{CA}$ 为围捕判据半径, $R_{\mathrm{safe}}$ 为软碰撞安全半径.

\begin{figure}[!htbp]
  \centering
  \includegraphics[width=0.75\linewidth]{参考资料/TASE_终稿提交版/problem.pdf}
  \bicaption{多无人艇协同围捕博弈场景示意}{Illustration of multi-USV cooperative hunting game}
  \label{fig:hunt_scene}
\end{figure}

\subsection{围捕成功数学判据}
\label{subsec:hunt_criterion}
为明确定义协同围捕的成功标准, 本文首先给出包围区域和逃逸者可达区域的概念.

\textbf{定义 4.1 \ (包围区域)} \quad 以第 $m$ 艘逃逸者 $E_m$ 为几何参考, 由 $N$
艘追击者及其运动状态张成的包围区域 $\Omega_m(t)$ 定义为
\begin{equation}
\Omega_m(t) = \left\{ (x,y)\ \Big|\ \|(x,y) - \mathbf{p}_m^C(t)\| \le d_{CA}\pm\sigma \right\},
\label{eq:besieged_area}
\end{equation}
其中 $\mathbf{p}_m^C(t)$ 为由追击者围绕 $E_m$ 所形成编队的几何中心, $\sigma$ 为
允许半径动态调节的容差因子.

\textbf{定义 4.2 \ (可达活动区域)} \quad 设逃逸者 $E_m$ 在任意策略作用下, 从当前
位置出发在未来 $\Delta t$ 时间内所能到达的所有位置点构成其可达活动区域 $R_e^m(t)$,
其大小取决于 $E_m$ 的速度限制与周围环境约束.

基于上述定义, 给出协同围捕成功的显式数学判据:

\textbf{判据 4.1 \ (协同围捕成功判据)} \quad 对于逃逸者 $E_m$, 若
\begin{equation}
\left\{
\begin{aligned}
\mathbf{p}_{E_m}(t) &\in \Omega_m(t), \quad \forall t\ge t_0, \\
R_e^m(t) &\subseteq \Omega_m(t), \quad \forall t\ge t_0,
\end{aligned}
\right.
\label{eq:hunt_cond}
\end{equation}
则认为追击者对 $E_m$ 的协同围捕成功. 此时, 逃逸者不仅当前位置位于包围区域内部,
且未来可能到达的所有位置也被包围区域完全包含, 即逃逸者彻底失去了机动逃脱的可能.

为进一步量化围捕性能, 定义最优围捕问题为以下 min-max 优化问题:
\begin{equation}
\min_{\pi_{P}}\ \max_{(x,y)\in R_e^m(t)}\ d\!\left((x,y),\, \partial\Omega_m(t)\right),
\label{eq:min_max}
\end{equation}
其中 $\pi_P$ 为追击者联合策略, $\partial\Omega_m$ 为包围区域的边界,
$d(\cdot,\cdot)$ 为点到集合的最短欧氏距离. 该优化问题的含义是: 追击者应设计策略,
使得逃逸者可达活动区域的最远点到包围边界的距离达到最小, 即最大程度压缩逃逸者的
机动空间.

\subsection{安全与非碰撞约束}
\label{subsec:safety_constraint}
为避免在围捕过程中发生碰撞, 所有智能体之间需满足软碰撞约束
\begin{equation}
d_{ij} \ge 2R_{\mathrm{safe}}, \quad \forall i\ne j,
\label{eq:collision}
\end{equation}
其中 $d_{ij}$ 为第 $i$ 艘与第 $j$ 艘无人艇的欧氏距离. 此外, 追击者需与障碍物保持
一定距离, 避免碰撞.

\subsection{基本假设}
\label{subsec:hunt_assumption}
为便于策略设计, 本章引入如下假设:

\textbf{假设 4.1} \quad 由于传感范围、通信带宽与计算能力的限制, 每艘追击者仅能获
取逃逸者的位置数据、距离其最近的两艘追击者状态以及最近三个障碍物信息.

\textbf{假设 4.2} \quad 所有追击者为同构无人艇, 其最大航速大于逃逸者最大航速, 即
$V_{\max}^E < V_{\max}^P$.

\textbf{注记 4.1} \quad 与部分传统研究仅将围捕视为 ``追击者到达指定包围区域'' 不
同, 本文的判据 4.1 通过 $R_e^m(t)\subseteq \Omega_m(t)$ 刻画了围捕对逃逸者可达
空间的实质约束, 更契合实际海战中包围压缩敌舰机动空间的战术意图.

\section{威胁势场与三阶段围捕判据}
\label{sec:tpf}

\subsection{威胁势场数学建模}
\label{subsec:tpf_model}
为刻画围捕过程中各智能体所处的态势, 本文以每艘逃逸者 $E_m$ 为中心, 构建一个以其
位置为原点的威胁势场 (Threat Potential Field, TPF). 对于博弈空间中任意一点
$\mathbf{q}=(x,y)$, 其威胁势能值定义为
\begin{equation}
U_{TPF}^{m}(\mathbf{q}) = -k_r\, e^{-\|\mathbf{q}-\mathbf{p}_{E_m}\|^2 / \sigma_r^2}
+ k_a\, \|\mathbf{q}-\mathbf{p}_{E_m}\|^2,
\label{eq:tpf}
\end{equation}
其中第一项为围绕 $E_m$ 形成的高斯型排斥势, 反映逃逸者附近对追击者的威胁;
第二项为远距离吸引势, 引导追击者从远处接近; $k_r, k_a, \sigma_r$ 为设计参数.
图~\ref{fig:tpf_field} 展示了威胁势场在博弈空间中的三维分布与俯视投影.

\begin{figure}[!htbp]
  \centering
  \subcaptionbox{三维能量分布\label{fig:tpf_3d}}{\includegraphics[width=0.45\linewidth]{参考资料/TASE_终稿提交版/blue_E.pdf}}
  \hfill
  \subcaptionbox{俯视等高线投影\label{fig:tpf_top}}{\includegraphics[width=0.45\linewidth]{参考资料/TASE_终稿提交版/asdads.pdf}}
  \bicaption{威胁势场的能量分布}{Threat potential field around the evaders}
  \label{fig:tpf_field}
\end{figure}

此外, 追击者的观测空间基于自身位置、最近两艘追击者及最近三个障碍物的相对信息构
造, 示意如图~\ref{fig:hunt_obs} 所示, 这一设计既保证了观测维度的可扩展性, 又有
效降低了冗余信息.

\begin{figure}[!htbp]
  \centering
  \includegraphics[width=0.72\linewidth]{参考资料/TASE_终稿提交版/position.pdf}
  \bicaption{协同围捕观测信息示意}{Schematic of the observational information for cooperative hunting}
  \label{fig:hunt_obs}
\end{figure}

\subsection{能量波峰 EC 与能量波谷 ET 的概念提出}
\label{subsec:ec_et}
传统 TPF 将势场值直接作为代价函数, 并未体现围捕过程的阶段性特征. 为解决此问题,
本文在 TPF 基础上引入\textbf{能量波峰 (Energy Crest, EC)}与\textbf{能量波谷
(Energy Trough, ET)}两个关键概念.

\textbf{定义 4.3 \ (能量波峰 EC)} \quad 逃逸者运动方向前方形成的威胁高能区域, 其
定义为
\begin{equation}
\mathrm{EC}_m(t) = \left\{ \mathbf{q}\ \Big|\ U_{TPF}^{m}(\mathbf{q}) \ge U_{EC},\ (\mathbf{q}-\mathbf{p}_{E_m})^T \mathbf{v}_{E_m} > 0 \right\},
\label{eq:ec}
\end{equation}
其中 $U_{EC}$ 为波峰阈值. 能量波峰是追击者希望避免进入的区域 (否则会被逃逸者直接
摆脱).

\textbf{定义 4.4 \ (能量波谷 ET)} \quad 逃逸者侧翼与后方的低能区域, 是追击者理想
的占位点, 定义为
\begin{equation}
\mathrm{ET}_m(t) = \left\{ \mathbf{q}\ \Big|\ U_{ET}^{\min} \le U_{TPF}^{m}(\mathbf{q}) \le U_{ET}^{\max},\ (\mathbf{q}-\mathbf{p}_{E_m})^T\mathbf{v}_{E_m} \le 0 \right\},
\label{eq:et}
\end{equation}
其中 $U_{ET}^{\min}, U_{ET}^{\max}$ 为波谷的势能界限. 能量波谷从几何上等价于逃
逸者 ``背后的安全扇区'', 是完成合围的最佳占位.

\subsection{围捕过程三阶段划分}
\label{subsec:three_stages}
基于 EC 与 ET 的定义, 本文将围捕过程划分为三个阶段:

\textbf{阶段 1 \ 接近阶段.} 追击者与逃逸者相距较远, 目标是迅速缩短距离. 此阶段奖
励函数侧重距离缩减分量.

\textbf{阶段 2 \ 成型阶段.} 追击者已接近逃逸者, 目标是快速占据各自对应的 ET 位
置, 形成初步合围. 此阶段奖励函数强化 ET 占位奖励与编队维持奖励.

\textbf{阶段 3 \ 保持阶段.} 追击者已占据 ET 位置, 目标是维持包围态势并压缩 $R_e^m(t)$
直到满足判据 4.1. 此阶段奖励函数强化包围完整性奖励.

阶段之间的切换由以下指标判断: 平均追-逃距离、追击者占据 ET 的比例、编队紧密度等.

\subsection{基于威胁势场的奖励机制}
\label{subsec:tpf_reward}
综合考虑三阶段围捕特点, 本章设计以下分阶段自适应奖励
\begin{equation}
r_t = \alpha_1(\kappa) r_{\mathrm{dist}} + \alpha_2(\kappa) r_{\mathrm{ET}} + \alpha_3(\kappa) r_{\mathrm{form}} + \alpha_4(\kappa) r_{\mathrm{obs}} + r_{\mathrm{term}},
\label{eq:hunt_reward}
\end{equation}
其中 $\kappa\in\{1,2,3\}$ 表示当前所处阶段, $\alpha_i(\kappa)$ 为各分量在不同阶
段的权重. 各分量含义如下:

\textbf{(1) 距离奖励 $r_{\mathrm{dist}}$:} 鼓励追击者缩短与目标逃逸者的距离, 尤
其在阶段 1 中占主导地位;

\textbf{(2) ET 占位奖励 $r_{\mathrm{ET}}$:} 当追击者进入 ET 区域时给予正奖励, 鼓
励占据理想位置, 在阶段 2、3 中权重较高;

\textbf{(3) 编队维持奖励 $r_{\mathrm{form}}$:} 奖励追击者形成均匀分布的包围编队,
避免扎堆;

\textbf{(4) 避障/避碰惩罚 $r_{\mathrm{obs}}$:} 对接近障碍物或其他追击者的行为给
予连续惩罚;

\textbf{(5) 终止奖励 $r_{\mathrm{term}}$:} 判据 4.1 成立时给予较大正奖励, 超时或
碰撞时给予较大负奖励.

\section{基于多头注意力的 MHD-MAPPO 算法}
\label{sec:mhd_mappo}

\subsection{MAPPO 算法框架}
\label{subsec:mappo_recap}
多智能体近端策略优化 (MAPPO) 算法 (参见 2.4.5 节) 为每个智能体配备独立的 Actor,
共享一个基于全局状态的 Critic. 在协同围捕场景中, 各 Actor 以局部观测 $o_i$ 为输
入输出动作 $a_i$, Critic 以全局联合状态 $\mathbf{s} = [o_1, o_2, \ldots, o_N]$
为输入输出状态值 $V_\phi(\mathbf{s})$, 协助 Actor 进行策略更新.

然而, 当智能体数量较多且观测含有大量冗余信息时 (例如远处的追击者和当前决策相关
性较弱), 传统 Critic 的全连接结构难以自适应地提取关键特征, 影响价值估计的准确
性. 为此, 本章在 Critic 中引入多头注意力机制.

\subsection{Critic 网络中的多头注意力机制}
\label{subsec:mha_critic}
设 Critic 网络接收到的全局输入为追击者观测序列 $\mathbf{X} = [o_1^T; o_2^T; \ldots; o_N^T]^T$.
首先通过三个线性变换将 $\mathbf{X}$ 映射为查询、键、值矩阵:
\begin{equation}
Q = \mathbf{X} W^Q,\quad K = \mathbf{X} W^K,\quad V = \mathbf{X} W^V,
\label{eq:qkv}
\end{equation}
随后按照式~\eqref{eq:attention} 与式~\eqref{eq:mha} 计算多头注意力输出
\begin{equation}
\mathbf{Z} = \mathrm{MHA}(Q, K, V) = \mathrm{Concat}(\mathrm{head}_1, \ldots, \mathrm{head}_h)\, W^O,
\end{equation}
之后经过层归一化、全连接与残差连接后输出状态值
\begin{equation}
V_\phi(\mathbf{s}) = f_{\mathrm{out}}\big(\mathrm{LayerNorm}(\mathbf{X} + \mathbf{Z})\big).
\label{eq:vc}
\end{equation}
与传统单头自注意力相比, 多头结构能够在多个子空间上并行建模智能体间、智能体-环境
间的多种相关性, 从而更全面地捕捉协同围捕中的关键态势.

\subsection{MHD-MAPPO 整体架构}
\label{subsec:mhd_mappo_arch}
MHD-MAPPO (Multi-Head attention Distributed MAPPO) 的整体架构如
图~\ref{fig:mhd_mappo} 所示. 每个追击者配备一个独立的 Actor 网络 $\pi_{\theta_i}(a_i|o_i)$,
Actor 网络采用全连接结构; Critic 网络则在全连接层前插入上述多头注意力模块, 用于
对全局状态进行特征提取. 训练阶段采用 CTDE 范式 (如图~\ref{fig:ctde} 所示), 执行
阶段各 Actor 仅依赖局部观测; Actor 与 Critic 网络结构详见图~\ref{fig:actor_critic}.

\begin{figure}[!htbp]
  \centering
  \includegraphics[width=0.95\linewidth]{参考资料/TASE_终稿提交版/model1.pdf}
  \bicaption{MHD-MAPPO 算法整体架构}{Overall architecture of MHD-MAPPO algorithm}
  \label{fig:mhd_mappo}
\end{figure}

\begin{figure}[!htbp]
  \centering
  \includegraphics[width=0.72\linewidth]{参考资料/TASE_终稿提交版/Drawing1.pdf}
  \bicaption{集中训练-分布式执行 (CTDE) 框架}{Centralized training and decentralized execution framework}
  \label{fig:ctde}
\end{figure}

\begin{figure}[!htbp]
  \centering
  \subcaptionbox{Actor 网络\label{fig:actor_net}}{\includegraphics[width=0.3\linewidth]{参考资料/TASE_终稿提交版/left.pdf}}
  \hfill
  \subcaptionbox{Critic 网络 (含 MHA)\label{fig:critic_net}}{\includegraphics[width=0.42\linewidth]{参考资料/TASE_终稿提交版/right.pdf}}
  \bicaption{MHD-MAPPO 中的 Actor 与 Critic 网络结构}{Actor and Critic network structures in MHD-MAPPO}
  \label{fig:actor_critic}
\end{figure}

\subsection{算法流程}
\label{subsec:mhd_mappo_algo}

\begin{algorithm}[!htbp]
  \caption{MHD-MAPPO 多无人艇协同围捕算法}
  \label{algo:mhd_mappo}
  \DontPrintSemicolon
  \SetKwInOut{Input}{输入}
  \SetKwInOut{Output}{输出}
  \Input{各 Actor 参数 $\{\theta_i\}_{i=1}^N$, 含 MHA 模块的 Critic 参数 $\phi$, 经验池 $\mathcal{D}$}
  \Output{训练完成的策略 $\{\pi_{\theta_i}\}_{i=1}^N$}
  初始化 Actor 与 MHA-Critic 网络参数\;
  \For{迭代 $k = 1, 2, \ldots, K$}{
    \For{并行环境 $e = 1, \ldots, E_p$}{
      每艘追击者 $P_i$ 利用 $\pi_{\theta_i^{\mathrm{old}}}$ 与环境交互, 采集轨迹\;
      根据当前所处阶段 $\kappa$ 按式~\eqref{eq:hunt_reward} 计算密集奖励 $r_t$\;
    }
    将所有轨迹存入经验池 $\mathcal{D}$\;
    使用 MHA-Critic 计算全局优势 $A_t$ 与目标值 $V_t^{\mathrm{target}}$\;
    \For{epoch $= 1, \ldots, E$}{
      按式~\eqref{eq:mappo_actor} 更新各 Actor 参数 $\theta_i$\;
      最小化均方误差更新 Critic 参数 $\phi$\;
    }
    同步 $\theta_i^{\mathrm{old}} \leftarrow \theta_i$\;
  }
\end{algorithm}

\section{基于人工势场的逃逸策略}
\label{sec:evader_strategy}
为训练追击者策略, 逃逸者需要具备相应的博弈对抗能力. 本章采用基于人工势场 (APF)
的逃逸策略: 追击者、障碍物以及作战水域边界均被建模为斥力源, 逃逸者在合成斥力作
用下选择逃逸方向. 对于逃逸者 $E_m$, 其受到的合力为
\begin{equation}
\mathbf{F}_{E_m} = -\nabla U_{\mathrm{tot}}^{m}(\mathbf{p}_{E_m}),
\end{equation}
其中
\begin{equation}
U_{\mathrm{tot}}^{m}(\mathbf{q}) = \sum_{n=1}^{N} U_{P_n}(\mathbf{q}) + \sum_{u=1}^{U} U_{\mathrm{obs}_u}(\mathbf{q}) + U_{\mathrm{bound}}(\mathbf{q}).
\end{equation}
逃逸者的艏向角与纵向速度根据合力方向与大小进行控制. 这一策略为追击者训练提供了合
理且具备挑战性的对抗对手.

\section{数值仿真实验与结果分析}
\label{sec:hunt_sim_exp}

\subsection{仿真环境与训练配置}
\label{subsec:hunt_config}
仿真环境设为 $200\,\mathrm{m}\times 200\,\mathrm{m}$ 的无界水域, 内部随机分布若
干圆形与多边形障碍物. 追击者数量 $N\in\{2,3,4\}$, 逃逸者数量 $M\in\{1,2\}$, 分别
构成 2v1、3v1、4v1、3v2、4v2 等多种博弈规模. 训练超参数与网络结构详见
表~\ref{tab:hunt_hyper}.

\begin{table}[htb]
  \centering
  \begin{minipage}[t]{0.9\linewidth}
    \caption{MHD-MAPPO 训练超参数}
    \label{tab:hunt_hyper}
    \begin{tabularx}{\linewidth}{lX}
      \toprule[1.5pt]
      {\heiti 超参数} & {\heiti 取值} \\\midrule[1pt]
      Actor/Critic 学习率 & $5\times 10^{-4}$ / $1\times 10^{-3}$ \\
      批大小 & $2048$ \\
      迭代总步数 & $3\times 10^{6}$ \\
      折扣因子 $\gamma$ & $0.99$ \\
      GAE 系数 $\lambda$ & $0.95$ \\
      裁剪范围 $\epsilon$ & $0.2$ \\
      注意力头数 $h$ & $4$ \\
      注意力特征维度 $d_k$ & $64$ \\
      并行环境数 & $16$ \\
      \bottomrule[1.5pt]
    \end{tabularx}
  \end{minipage}
\end{table}

\subsection{训练收敛性分析}
\label{subsec:hunt_train_curve}
图~\ref{fig:hunt_train_curve} 比较了 MHD-MAPPO 与基于单头自注意力的 MAPPO 变体
在 3v1 协同围捕场景下的训练曲线, 具体包括平均奖励、价值函数损失、策略损失与策略
熵四个指标. 可以看出, MHD-MAPPO 在各项指标上均展现出更快的收敛速度与更高的稳定
性: 平均奖励更早达到高水平并保持稳定, 价值损失与策略损失更快收敛至低值, 策略熵
衰减更平滑. 这充分证明了多头注意力机制对多智能体协同决策的有效性.

\begin{figure}[!htbp]
  \centering
  \subcaptionbox{平均奖励\label{fig:hunt_reward}}{\includegraphics[width=0.44\linewidth]{参考资料/TASE_终稿提交版/reward.pdf}}
  \hfill
  \subcaptionbox{价值损失\label{fig:hunt_vloss}}{\includegraphics[width=0.44\linewidth]{参考资料/TASE_终稿提交版/value_loss.pdf}}

  \vspace{8pt}
  \subcaptionbox{策略损失\label{fig:hunt_ploss}}{\includegraphics[width=0.44\linewidth]{参考资料/TASE_终稿提交版/policy_loss.pdf}}
  \hfill
  \subcaptionbox{策略熵损失\label{fig:hunt_ent}}{\includegraphics[width=0.44\linewidth]{参考资料/TASE_终稿提交版/entropy.pdf}}
  \bicaption{MHD-MAPPO 与基线算法训练曲线对比}{Training curves of MHD-MAPPO vs. baselines}
  \label{fig:hunt_train_curve}
\end{figure}

\subsection{典型场景仿真结果}
\label{subsec:hunt_scenarios}
图~\ref{fig:hunt_compare} 展示了在相同初始条件下, 确定性算法、启发式算法与本文
MHD-MAPPO 算法的协同围捕轨迹对比以及三艘追击者与逃逸者的距离曲线. 结果显示,
MHD-MAPPO 的围捕过程更为紧凑、协同性更强, 追-逃距离缩小更快. 图~\ref{fig:hunt_snapshots}
进一步展示了 MHD-MAPPO 在 $t=20$ s、$t=50$ s、$t=80$ s 三个典型时刻的博弈态势
快照, 可清晰观察到追击者依次进入接近阶段、成型阶段与保持阶段, 最终将逃逸者完全
包围于 $\Omega_m(t)$ 内.

\begin{figure}[!htbp]
  \centering
  \subcaptionbox{确定性算法\label{fig:hunt_det}}{\includegraphics[width=0.45\linewidth]{参考资料/TASE_终稿提交版/path3.eps}}
  \hfill
  \subcaptionbox{启发式算法\label{fig:hunt_heu}}{\includegraphics[width=0.45\linewidth]{参考资料/TASE_终稿提交版/path2.eps}}

  \vspace{6pt}
  \subcaptionbox{本文 MHD-MAPPO\label{fig:hunt_ours}}{\includegraphics[width=0.45\linewidth]{参考资料/TASE_终稿提交版/path1.eps}}
  \hfill
  \subcaptionbox{追逃距离曲线\label{fig:hunt_dist}}{\includegraphics[width=0.45\linewidth]{参考资料/TASE_终稿提交版/path1_data.eps}}
  \bicaption{三种协同围捕算法的轨迹与距离曲线对比}{Trajectory and distance curves of three cooperative hunting schemes}
  \label{fig:hunt_compare}
\end{figure}

\begin{figure}[!htbp]
  \centering
  \subcaptionbox{$t=20$ s (接近阶段)\label{fig:hunt_t20}}{\includegraphics[width=0.31\linewidth]{参考资料/TASE_终稿提交版/pic1.png}}
  \hfill
  \subcaptionbox{$t=50$ s (成型阶段)\label{fig:hunt_t50}}{\includegraphics[width=0.31\linewidth]{参考资料/TASE_终稿提交版/pic2.png}}
  \hfill
  \subcaptionbox{$t=80$ s (保持阶段)\label{fig:hunt_t80}}{\includegraphics[width=0.31\linewidth]{参考资料/TASE_终稿提交版/pic3.png}}
  \bicaption{MHD-MAPPO 协同围捕过程不同时刻快照}{Snapshots of MHD-MAPPO cooperative hunting at different time instants}
  \label{fig:hunt_snapshots}
\end{figure}

\subsection{含障碍环境下的泛化性实验}
\label{subsec:hunt_generalization}
为验证算法的泛化能力, 在训练集以外的障碍物分布、无人艇初始位置与逃逸者速度组合下
进行测试, 每种配置运行 100 次独立实验, 成功率与平均围捕时间如
表~\ref{tab:hunt_generalization} 所示. 可以看出, 所提算法在未见过的测试环境下仍
能保持 85\% 以上的围捕成功率, 泛化能力强.

\begin{table}[htb]
  \centering
  \begin{minipage}[t]{0.9\linewidth}
    \caption{不同测试配置下的泛化性实验结果}
    \label{tab:hunt_generalization}
    \begin{tabularx}{\linewidth}{lXXX}
      \toprule[1.5pt]
      {\heiti 测试配置} & {\heiti 成功率} & {\heiti 平均围捕时间} & {\heiti 平均路径长度} \\\midrule[1pt]
      训练集分布 & 92\% & 46.3 s & 148.7 m \\
      新障碍分布 & 88\% & 51.2 s & 161.4 m \\
      新初始位置 & 90\% & 49.8 s & 155.9 m \\
      快速逃逸者 & 85\% & 58.7 s & 178.3 m \\
      \bottomrule[1.5pt]
    \end{tabularx}
  \end{minipage}
\end{table}

\subsection{消融实验}
\label{subsec:hunt_ablation}
为验证多头注意力与阶段自适应奖励机制各自的贡献, 开展消融实验. 分别构造以下变体:
(a) 去除 MHA 模块 (保留单头注意力); (b) 去除阶段自适应, 使用固定奖励权重; (c) 去
除 ET 占位奖励; (d) 同时去除 MHA 与阶段自适应. 实验结果如
表~\ref{tab:hunt_ablation} 所示. 可以看出, 完整的 MHD-MAPPO 相比于各消融变体在
成功率与围捕时间上均取得了明显优势.

\begin{table}[htb]
  \centering
  \begin{minipage}[t]{0.9\linewidth}
    \caption{协同围捕消融实验结果}
    \label{tab:hunt_ablation}
    \begin{tabularx}{\linewidth}{lXX}
      \toprule[1.5pt]
      {\heiti 模型变体} & {\heiti 成功率} & {\heiti 平均围捕时间} \\\midrule[1pt]
      完整 MHD-MAPPO & \textbf{92\%} & \textbf{46.3} s \\
      (a) 去除 MHA & 83\% & 58.1 s \\
      (b) 去除阶段自适应 & 79\% & 63.7 s \\
      (c) 去除 ET 占位奖励 & 81\% & 60.5 s \\
      (d) 同时去除 MHA 与阶段自适应 & 68\% & 79.2 s \\
      \bottomrule[1.5pt]
    \end{tabularx}
  \end{minipage}
\end{table}

\section{多无人艇物理实验平台验证}
\label{sec:hardware_exp}

除数值仿真外, 本章进一步搭建多无人艇物理实验平台, 在真实水面环境下开展 2 对 1 与
3 对 1 协同围捕实物实验, 验证所提算法的工程可行性.

\subsection{多无人艇硬件平台}
\label{subsec:usv_hardware}
多无人艇物理实验平台由若干艘自主研制的小型水面无人艇组成, 每艘无人艇主要由船体、
推进系统、传感器模块、通信模块与机载计算单元构成, 整体参数见
表~\ref{tab:usv_hardware_spec}. 无人艇实物如图~\ref{fig:usv_photo} 所示.

\begin{table}[htb]
  \centering
  \begin{minipage}[t]{0.95\linewidth}
    \caption{多无人艇硬件平台主要参数}
    \label{tab:usv_hardware_spec}
    \begin{tabularx}{\linewidth}{lX}
      \toprule[1.5pt]
      {\heiti 组件/项目} & {\heiti 配置/说明} \\\midrule[1pt]
      船体长度 $\times$ 宽度 & $0.85$ m $\times$ $0.42$ m \\
      排水量 & 约 12 kg \\
      推进方式 & 双无刷电机差速推进 \\
      最大航速 & $1.8$ m/s \\
      动力电池 & 22.2 V 锂聚合物电池, 续航约 2 h \\
      惯性测量单元 (IMU) & MPU-9250, 9 轴 \\
      GNSS 定位 & RTK 厘米级 GPS 模块 \\
      机载计算单元 & NVIDIA Jetson Nano / Xavier NX \\
      通信链路 & 2.4 GHz Wi-Fi Mesh + 433 MHz 数传电台 \\
      \bottomrule[1.5pt]
    \end{tabularx}
  \end{minipage}
\end{table}

\begin{figure}[!htbp]
  \centering
  \includegraphics[width=0.75\linewidth]{参考资料/TASE_终稿提交版/exper.pdf}
  \bicaption{多无人艇实验平台整体结构与实物}{Structure and photograph of the multi-USV experimental platform}
  \label{fig:usv_photo}
\end{figure}

\subsection{通信与控制系统}
\label{subsec:comm_ctrl}
实验系统采用 ``岸基控制站 + 机载自主决策'' 的分布式架构. 岸基地面控制站 (Ground
Control Station, GCS) 负责任务下发、状态监控与数据记录; 各无人艇机载计算单元上
部署训练完成的 MHD-MAPPO 策略网络, 独立计算控制指令; 无人艇之间通过 Wi-Fi Mesh
广播自身状态, 实现必要的邻居信息交互. 控制层分两级: 上层 RL 策略给出期望的纵向
速度与艏向角速度, 下层通过 PID 控制器将指令映射为左右电机推力.

\subsection{仿真到实物的策略迁移 (Sim2Real)}
\label{subsec:sim2real}
从仿真环境到物理实验直接部署强化学习策略通常会遇到 ``Sim2Real'' 鸿沟, 具体体现在
质量分布、推进响应、传感器噪声等多方面的差异. 本章采用以下策略弥合这一鸿沟:

\textbf{(1) 领域随机化 (Domain Randomization).} 在训练阶段随机采样无人艇质量、推
进力矩响应、传感器测量噪声及海流速度等参数, 使策略对参数变化具有鲁棒性;

\textbf{(2) 动作延迟补偿.} 将仿真中的控制指令施加一定随机延迟 ($50\sim 200$ ms),
模拟真实通信与控制链路的时滞;

\textbf{(3) 观测噪声注入.} 对仿真环境中的位置、姿态观测添加与真实传感器相当的白
噪声与偏差漂移;

\textbf{(4) 微调训练.} 在物理平台上采集少量实验数据, 对策略网络进行少量迭代的微
调, 进一步提升实物性能.

\subsection{实验场景与方案}
\label{subsec:hunt_exp_scene}
实验在上海大学无人艇试验水池与附近湖面开展. 实验水域尺寸约
$30\,\mathrm{m}\times 25\,\mathrm{m}$, 四周由安全警戒浮标围成, 以防无人艇冲出测
试区域. 每组实验运行 10 次, 记录各无人艇的 GPS 轨迹、速度与艏向数据, 并同步摄影
存档.

\subsection{实物实验运动状态}
\label{subsec:exp_motion}
图~\ref{fig:hunt_exp_state} 展示了物理实验中各追击者的纵向速度与艏向角速度随时
间变化曲线. 可以观察到, 在接近阶段 (约 0$\sim$15 s), 追击者以较高速度快速逼近逃
逸者, 艏向频繁调整以修正追击方向; 在成型阶段 (约 15$\sim$30 s), 追击者速度略有
降低, 艏向趋于稳定, 占据 ET 位置; 在保持阶段 (约 30 s 之后), 追击者速度与艏向均
趋于平稳, 形成稳定的合围编队. 上述运动数据与仿真结果一致, 体现了 Sim2Real 迁移
后的策略可靠性.

\begin{figure}[!htbp]
  \centering
  \subcaptionbox{追击者纵向速度\label{fig:hunt_v}}{\includegraphics[width=0.8\linewidth]{参考资料/TASE_终稿提交版/v.eps}}

  \vspace{4pt}
  \subcaptionbox{追击者艏向角速度\label{fig:hunt_yaw}}{\includegraphics[width=0.8\linewidth]{参考资料/TASE_终稿提交版/yawspeed.eps}}
  \bicaption{物理实验中追击者的运动状态}{Motion states of pursuers in physical experiments}
  \label{fig:hunt_exp_state}
\end{figure}

\subsection{2 对 1 协同围捕实验结果}
\label{subsec:2v1_exp}
在 2 对 1 实物实验中, 两艘追击者分别从左侧与右侧接近逃逸者, 在 $t\approx 18$ s
时开始进入成型阶段, 分别占据逃逸者左右 ET 位置; 在 $t\approx 35$ s 时满足判据 4.1,
围捕成功. 10 次重复实验的平均围捕时间为 $37.8$ s, 成功率为 $90\%$.

\subsection{3 对 1 协同围捕实验结果}
\label{subsec:3v1_exp}
在 3 对 1 实物实验中, 三艘追击者呈三角形分布将逃逸者包围在中心区域, 从
$t\approx 22$ s 开始稳定维持包围编队直至捕获. 由于多艇协同效果更好, 3 对 1 实验的
平均围捕时间仅为 $29.5$ s, 成功率达 $95\%$, 优于 2 对 1 场景.

\subsection{实物实验鲁棒性分析}
\label{subsec:hunt_exp_robust}
表~\ref{tab:hunt_exp_summary} 总结了不同实验条件下的成功率与围捕时间. 可以看出,
即便在轻微风浪 (水面波纹高度 $\le 5$ cm、表面流速 $\le 0.1$ m/s) 等扰动条件下,
所提算法仍能保持较高的围捕成功率, 充分验证了 MHD-MAPPO 算法在实际海上环境中的
工程可行性.

\begin{table}[htb]
  \centering
  \begin{minipage}[t]{0.95\linewidth}
    \caption{不同实验条件下物理实验结果汇总}
    \label{tab:hunt_exp_summary}
    \begin{tabularx}{\linewidth}{lXXXX}
      \toprule[1.5pt]
      {\heiti 实验条件} & {\heiti 场景} & {\heiti 成功率} & {\heiti 平均时间} & {\heiti 平均距离误差} \\\midrule[1pt]
      水池无扰动 & 2v1 & 90\% & 37.8 s & 0.31 m \\
      水池无扰动 & 3v1 & 95\% & 29.5 s & 0.28 m \\
      湖面轻微风浪 & 2v1 & 85\% & 42.4 s & 0.45 m \\
      湖面轻微风浪 & 3v1 & 90\% & 33.7 s & 0.39 m \\
      \bottomrule[1.5pt]
    \end{tabularx}
  \end{minipage}
\end{table}

\section{本章小结}
\label{sec:hunt_summary}
本章针对复杂环境下多无人艇协同围捕博弈问题开展了系统研究: 首先, 基于六自由度非线
性运动学模型, 给出了协同围捕成功的显式数学判据, 填补了现有研究在评价标准上的空
白; 其次, 提出能量波峰 EC 与能量波谷 ET 概念, 将围捕过程划分为接近-成型-保持三阶
段, 设计了阶段自适应的密集奖励机制; 再次, 在 MAPPO 的 Critic 网络中引入多头注意
力机制, 构建了 MHD-MAPPO 算法, 显著提升了多艇协同决策效率; 最后, 搭建了多无人艇
物理实验平台, 在真实水面环境下完成 2 对 1 与 3 对 1 协同围捕实物实验, 验证算法从
仿真到实物迁移的工程可行性, 是本章研究的关键创新点之一. 本章工作为多无人艇在现实
海洋环境中的博弈对抗奠定了算法与工程基础.
